\documentclass[11pt,a4paper]{article}
\usepackage[utf8]{inputenc}
\usepackage[T1]{fontenc}
\usepackage{amsmath, amssymb}
\usepackage{amsthm}
\usepackage[margin=2.5cm]{geometry}
\usepackage{hyperref}
\hypersetup{pdftitle={Energy Decay in a Damped Harmonic Oscillator}}

\newtheorem{theorem}{Theorem}

\title{Energy Decay in a Damped Harmonic Oscillator}
\date{October 2026}

\begin{document}
\maketitle

\begin{abstract}
The damped harmonic oscillator is the simplest model of a system that loses energy to its surroundings [1]. Its equation of motion is
\end{abstract}

\section{Introduction}\label{sec:intro}

Write straight onto the page. Press \textbf{/} on an empty line for sections, figures and tables, and \textbf{\$} anywhere in a sentence for a formula such as $\sigma_{x}^{2}=\langle x^{2}\rangle-\langle x\rangle^{2}$ --- click it to edit it again. The LaTeX on the right is written for you as you type.
\begin{equation}
  \int_{-\infty}^{\infty}e^{-x^{2}}\,dx=\sqrt{\pi}
  \label{eq:gauss}
\end{equation}
m x'' + c x' + k x = 0

where m \textgreater{} 0 is the mass, c \textgreater{}= 0 the damping coefficient and k \textgreater{} 0 the spring constant. [footnote]All quantities are in SI units.[/footnote] Dividing by m and writing omega\_0 = sqrt(k/m) and gamma = c/(2m) gives the standard form
\begin{equation}
  
\end{equation}

\begin{equation}
  
\end{equation}

\section{Solution}

Substituting x = e\textasciicircum{}(lambda*t) into equation (2) gives the characteristic equation lambda\textasciicircum{}2 + 2*gamma*lambda + omega\_0\textasciicircum{}2 = 0, whose roots are
\begin{equation}
  
\end{equation}

\begin{theorem}
  Differentiate E and use equation (1): dE/dt = m*x'*x'' + k*x*x' = x'*(m*x'' + k*x) = -c*x'\textasciicircum{}2.
\end{theorem}

\begin{proof}
  Numerical check
\end{proof}

We integrated equation (2) with omega\_0 = 2, gamma = 0.2 and x\_0 = 1 on t in [0, 10], using the explicit Euler method and the classical fourth-order Runge-Kutta method (RK4). For each step size h we: 1. integrated from t = 0 to t = 10 with 10/h steps; 2. compared every step with the exact solution; 3. recorded the largest absolute difference.

\section{}

\end{document}
